\quad O Merge Sort é um algoritmo de ordenação baseado na estratégia
de divisão e conquista. Seu funcionamento consiste em dividir o vetor
em partes menores, ordenar essas partes recursivamente e, em seguida,
combinar os resultados para formar um único vetor ordenado.

\quad Inicialmente, o vetor é dividido em duas metades. Esse processo
continua até que cada parte contenha apenas um elemento, situação em
que ela já está ordenada. Em seguida, a função \textit{merge} combina
as duas metades, comparando seus elementos e colocando-os em ordem
crescente no vetor original.

\quad A função \textit{mergeSort} é responsável pelas divisões
recursivas. Ela calcula o índice do meio do vetor, ordena a metade
esquerda, ordena a metade direita e chama a função \textit{merge}
para intercalar os resultados. A função \textit{merge}, por sua vez,
utiliza dois vetores auxiliares para armazenar temporariamente os
elementos das duas metades durante a intercalação.

\begin{lstlisting}[caption={Implementação do Merge Sort}]
void merge(int* vetor, int esq, int meio, int dir) {
    int n1 = meio - esq + 1;
    int n2 = dir - meio;

    int* E = new int[n1];
    int* D = new int[n2];

    for (int i = 0; i < n1; i++)
        E[i] = vetor[esq + i];

    for (int j = 0; j < n2; j++)
        D[j] = vetor[meio + 1 + j];

    int i = 0, j = 0;
    int k = esq;

    while (i < n1 && j < n2) {
        if (E[i] <= D[j]) {
            vetor[k] = E[i];
            i++;
        } else {
            vetor[k] = D[j];
            j++;
        }
        k++;
    }

    while (i < n1) {
        vetor[k] = E[i];
        i++;
        k++;
    }

    while (j < n2) {
        vetor[k] = D[j];
        j++;
        k++;
    }

    delete[] E;
    delete[] D;
}

void mergeSort(int* vetor, int tamanho,
               int esq = 0, int dir = -1) {
    if (dir == -1)
        dir = tamanho - 1;

    if (esq >= dir)
        return;

    int meio = esq + (dir - esq) / 2;

    mergeSort(vetor, tamanho, esq, meio);
    mergeSort(vetor, tamanho, meio + 1, dir);

    merge(vetor, esq, meio, dir);
}
\end{lstlisting}

\quad Ao utilizar a função \textit{mergeSort}, basta informar o vetor
e seu tamanho. Os parâmetros \textit{esq} e \textit{dir} possuem
valores padrão, permitindo que a primeira chamada considere todo o
vetor. Nas chamadas recursivas, esses parâmetros delimitam as partes
que serão ordenadas.

\quad O processo de divisão termina quando o intervalo possui zero
ou um elemento, condição identificada por \textit{esq} maior ou igual
a \textit{dir}. Após as chamadas recursivas, as partes ordenadas são
intercaladas até que todo o vetor esteja organizado em ordem crescente.
